% Part: first-order-logic
% Chapter: syntax-and-semantics
% Section: unique-readability

\documentclass[../../../include/open-logic-section]{subfiles}

\begin{document}

\olfileid{fol}{syn}{unq}

\olsection{Mung Bisa Diwaca Kanthi Siji Cara}

\begin{explain}
Carane kita nemtokake !!{formula}s njamin yen saben !!{formula} mung
bisa \emph{diwaca kanthi siji cara}, yaiku miturut aturan pambentukan
!!{formula}s, saktemene mung ana siji cara kanggo mbangun lan
nafsirake formula kasebut. Yen ora mangkono, bakal ana !!{formula}s
kang ambigu, yaiku !!{formula}s kang bisa diwaca utawa ditafsirake
luwih saka siji cara---lan cetha yen kita kepengin nyingkiri perkara
iku. Nanging, kang luwih wigati, tanpa sipat iki umume definisi lan
bukti kang bakal kita wenehake ora bakal lumaku.

Mbokmenawa cara paling cetha kanggo nerangake iki yaiku ndeleng apa
kang bakal kelakon yen kita menehi aturan pambentukan !!{formula}s
kang ora njamin yen formula mung bisa diwaca kanthi siji cara.
Umpamane, kita bisa wae lali nyertakake kurung ing aturan pambentukan
kanggo konektif, kayata ngidini aturan iki:
\begin{quote}
Yen $!A$ lan $!B$ iku !!{formula}s, mula $!A \lif !B$ uga formula.
\end{quote}
Miwiti saka formula atomik $!D$, aturan iki ngidini kita mbangun
$!D \lif !D$. Saka iki bebarengan karo $!D$, kita bakal oleh $!D \lif !D
\lif !D$. Nanging, ana rong cara kanggo nindakake iki:
\begin{enumerate}
\item Kita nganggep $!D$ minangka $!A$ lan $!D \lif !D$ minangka $!B$.
\item Kita nganggep $!A$ minangka $!D \lif !D$ lan $!B$ yaiku~$!D$.
\end{enumerate}
Mula, ana rong cara kanggo maca !!{formula}~$!D \lif !D \lif !D$.
Formula iku awujud $!B \lif !C$ kanthi $!B$ yaiku $!D$ lan $!C$
yaiku $!D \lif !D$, nanging \emph{uga} awujud $!B \lif !C$ kanthi
$!B$ yaiku $!D \lif !D$ lan $!C$ yaiku~$!D$.

Yen iki kelakon, definisi kita ora mesthi lumaku. Umpamane, nalika
kita nemtokake !!{main operator} sawijining formula, kita kandha:
ing formula kang awujud $!B \lif !C$, !!{main operator} yaiku
kemunculan~$\lif$ kang dituduhake. Nanging yen kita bisa nyocogake
formula $!D \lif !D \lif !D$ karo $!B \lif !C$ kanthi rong cara mau,
ing kasus kapisan kemunculan $\lif$ kang kapisan dadi !!{main operator},
dene ing kasus kapindho kemunculan kang kapindho. Nanging kita
ngersakake supaya !!{main operator} dadi \emph{fungsi} saka
!!{formula}, yaiku saben !!{formula} kudu duwe pas siji kemunculan
!!{main
  operator}.
\end{explain}

\begin{lem}
Cacahing kurung buka lan kurung tutup ing !!a{formula}~$!A$ iku padha.
\end{lem}

\begin{proof}
Kita mbuktekake iki kanthi induksi miturut cara $!A$ kawangun. Iki
mbutuhake rong bab: (a) Kaping pisanan, kita kudu mbuktekake yen kabeh
formula atomik nduweni sipat kang dikarepake (dhasar induksi).
(b) Banjur, kita kudu mbuktekake yen nalika kita mbangun formula anyar
saka formula kang wis ana, formula anyar nduweni sipat kasebut yen
formula lawas uga nduweni.

Upamana $l(!A)$ iku cacahing kurung buka, lan $r(!A)$ cacahing
kurung tutup ing~$!A$; mangkono uga $l(t)$ lan $r(t)$ kanggo cacahing
kurung buka lan kurung tutup ing sawijining term~$t$.

\begin{prob}
    Buktèkna yen kanggo saben term~$t$, $l(t) = r(t)$.
\end{prob}

\begin{enumerate}
\tagitem{prvFalse}{\indcase{!A}{\lfalse}{$\indfrm$ nduweni $0$ kurung
    buka lan $0$ kurung tutup.}}{}

\tagitem{prvTrue}{\indcase{!A}{\ltrue}{$\indfrm$ nduweni $0$ kurung
    buka lan $0$ kurung tutup.}}{}

\item \indcase{!A}{\Atom{R}{t_1,\dots,t_n}}{$l(\indfrm) = 1 + l(t_1) +
  \dots + l(t_n) = 1 + r(t_1) + \dots + r(t_n) = r(\indfrm)$. Ing kene
  kita nggunakake kasunyatan, kang ditinggal minangka latihan, yen
  $l(t) = r(t)$ kanggo saben term~$t$.}

\item \indcase{!A}{\eq[t_1][t_2]}{$l(\indfrm) = l(t_1) + l(t_2) =
  r(t_1) + r(t_2) = r(\indfrm)$.}

\tagitem{prvNot}{\indcase{!A}{\lnot !B}{Miturut hipotesis induksi,
    $l(!B) = r(!B)$. Mula $l(\indfrm) = l(!B) = r(!B) =
    r(\indfrm)$.}}{}

\item \indcase{!A}{(!B \ast !C)}{Miturut hipotesis induksi, $l(!B) =
  r(!B)$ lan $l(!C) = r(!C)$. Mula $l(\indfrm) = 1 + l(!B) + l(!C) =
  1 + r(!B) + r(!C) = r(\indfrm)$.}

\tagitem{prvAll}{\indcase{!A}{\lforall[x][!B]}{Miturut hipotesis
    induksi, $l(!B) = r(!B)$. Mula, $l(\indfrm) = l(!B) = r(!B) =
    r(\indfrm)$.}}{}

% Print case for \lexists if prvEx and not prvAll
\tagitem{prvAll,notprvEx}{}{\indcase{!A}{\lexists[x][!B]}{Miturut hipotesis
    induksi, $l(!B) = r(!B)$. Mula, $l(\indfrm) = l(!B) = r(!B) =
    r(\indfrm)$.}}

% Just say similarly if prvEx and prvAll
\tagitem{notprvAll,notprvEx}{}{\indcase{!A}{\lexists[x][!B]}{Semono uga.}}
\end{enumerate}
\end{proof}

\begin{defn}[Awalan sejati]
Sawijining rerangken simbol $!B$ iku \emph{awalan sejati} saka
rerangken simbol~$!A$ yen nyambungake $!B$ karo rerangken simbol
kang ora kosong ngasilake~$!A$.
\end{defn}

\begin{lem}\ollabel{lem:no-prefix}
Yen $!A$ iku !!a{formula}, lan $!B$ iku awalan sejati saka $!A$,
mula $!B$ dudu !!a{formula}.
\end{lem}

\begin{proof}
Latihan.
\end{proof}

\begin{prob}
Buktèkna \olref[fol][syn][unq]{lem:no-prefix}.
\end{prob}

\begin{prop}
\ollabel{prop:unique-atomic}
Yen $!A$ iku !!{formula} atomik, mula formula kasebut nyukupi siji, lan mung
siji, saka syarat-syarat ing ngisor iki.
\begin{enumerate}
\tagitem{prvFalse}{$!A \ident \lfalse$.}{}
\tagitem{prvTrue}{$!A \ident \ltrue$.}{}
\item $!A \ident \Atom{R}{t_1,\dots,t_n}$ kanthi $R$ iku
  !!{predicate} $n$-panggonan, $t_1$, \dots, $t_n$ iku term, lan
  saben $R$, $t_1$, \dots, $t_n$ ditemtokake kanthi siji cara.
\item $!A \ident \eq[t_1][t_2]$ kanthi $t_1$ lan $t_2$ iku term
  kang ditemtokake kanthi siji cara.
\end{enumerate}
\end{prop}

\begin{proof}
Latihan.
\end{proof}

\begin{prob}
Buktèkna \olref[fol][syn][unq]{prop:unique-atomic} (Pituduh:
Rumusna lan buktèkna versi \olref[fol][syn][unq]{lem:no-prefix}
kanggo term.)
\end{prob}

\begin{prop}[Mung Bisa Diwaca Kanthi Siji Cara]
Saben !!{formula} nyukupi siji, lan mung siji, saka syarat-syarat
ing ngisor iki.
\begin{enumerate}
\item $!A$ iku atomik.

\tagitem{prvNot}{$!A$ awujud $\lnot !B$.}{}

\tagitem{prvAnd}{$!A$ awujud $(!B \land !C)$.}{}

\tagitem{prvOr}{$!A$ awujud $(!B \lor !C)$.}{}

\tagitem{prvIf}{$!A$ awujud $(!B \lif !C)$.}{}

\tagitem{prvIff}{$!A$ awujud $(!B \liff !C)$.}{}

\tagitem{prvAll}{$!A$ awujud $\lforall[x][!B]$.}{}

\tagitem{prvEx}{$!A$ awujud $\lexists[x][!B]$.}{}
\end{enumerate}
Kajaba iku, ing saben kasus $!B$, utawa $!B$ lan $!C$, ditemtokake
kanthi siji cara. Iki ateges, umpamane, ora ana pasangan beda $!B$,
$!C$ lan $!B'$, $!C'$ saengga $!A$ awujud loro-lorone
\iftag{prvIf}{$(!B \lif !C)$ lan $(!B' \lif !C')$.}{
    \iftag{prvOr}{$(!B \lor !C)$ lan $(!B' \lor !C')$.}{
      \iftag{prvAnd}{$(!B \land !C)$ lan $(!B' \land !C')$.}{}}}
\end{prop}

\begin{proof}
Aturan pambentukan mbutuhake yen !!a{formula} ora atomik, formula iku
kudu diwiwiti tandha kurung buka~(, \iftag{prvNot}{$\lnot$,}{}
utawa kuantor. Kosok baline, saben !!{formula} kang diwiwiti salah
siji saka simbol-simbol iki mesthi atomik: !!a{predicate},
!!a{function}, !!a{constant}\iftag{prvFalse}{, $\lfalse$}{}\iftag{prvTrue}{, $\ltrue$}{}.

Dadi, kita mung perlu nuduhake yen $!A$ awujud $(!B \ast
!C)$ lan uga awujud $(!B' \mathbin{\ast'} !C')$, mula $!B \ident
!B'$, $!C \ident !C'$, lan $\ast = {\ast'}$.

Upamana $!A \ident (!B \ast !C)$ lan $!A \ident (!B'
\mathbin{\ast'} !C')$. Mula bisa uga $!B \ident !B'$ utawa ora.
Yen padha, cetha yen $\ast = {\ast'}$ lan $!C \ident !C'$,
amarga loro-lorone banjur dadi subrerangken saka $!A$ kang diwiwiti
ing papan kang padha lan dawane padha. Kasus liyane yaiku
$!B \not\ident !B'$. Amarga $!B$ lan $!B'$ loro-lorone subrerangken
saka $!A$ kang diwiwiti ing papan kang padha, salah siji mesthi
dadi awalan sejati saka sijine. Nanging iki ora mungkin miturut
\olref{lem:no-prefix}.
\end{proof}


\end{document}
